\section{Measuring Power of Star-Free Languages}\label{sec:sf}

In this section, we give a proof of the first main theorem, which gives a decidable characterisation of \(\sf\)-measurability for regular languages.
The key lemma of the proof of Theorem~\ref{thm:sf} is the following lemma.
\begin{lemma}\label{lem:cat}
 The class \(\ext(\gd)\) is closed under concatenation.
\end{lemma}

\begin{theorem}\label{thm:sf}
The class of star-free languages and the class of generalised definite languages have the same measuring power, \ie, \(\ext(\sf) = \ext(\gd)\).
\end{theorem}
\begin{proof}[Proof of Theorem~\ref{thm:sf}]
By definition, \(\fin\) is contained in \(\gd\).
Moreover, \(\ext(\gd)\) is closed under Boolean operations by Lemma~\ref{lem:closure} since \(\gd\) is so.
Combining these facts with Lemma~\ref{lem:cat}, we can conclude that
\(\sf \subseteq \ext(\gd)\), which implies 
\(\ext(\sf) \subseteq \ext(\ext(\gd)) = \ext(\gd) \subseteq \ext(\sf)\)
by Lemma~\ref{lem:closure}.
That is, \(\ext(\sf) = \ext(\gd)\) holds.
\end{proof}
\begin{corollary}
 The \(\sf\)-measurability for a given deterministic automaton is decidable in quadratic time.
\end{corollary}

The proof of Lemma~\ref{lem:cat} requires some additional notion and lemma.
For a given generalised definite language \(L = F A^* G\) where \(F, G\) are finite, we define its \emph{prefix density}
\(\pred(L)\) as \(\d(F A^*)\) and its
\emph{suffix density} \(\sufd(L)\) as \(\d(A^* G)\).
By Lemma~\ref{lem:density}, we have \(\d(L) = \pred(L) \cdot \sufd(L)\).